\begin{thebibliography}{99}
\bibitem{alphaproof-ref1} de Moura, L. \& Ullrich, S. The Lean 4 theorem prover and programming language. In Automated Deduction--CADE 28: Proc. 28th International Conference on Automated Deduction Proceedings (eds Platzer, A. \& Sutcliffe, G.) 625--635 (Springer, 2021). 题名中译: Lean 4 定理证明器与编程语言.
\bibitem{alphaproof-ref2} Silver, D. et al. A general reinforcement learning algorithm that masters chess, shogi, and Go through self-play. Science 362, 1140--1144 (2018). 题名中译: 通过自我博弈掌握国际象棋、将棋和围棋的通用强化学习算法.
\bibitem{alphaproof-ref3} Chervonyi, Y. et al. Gold-medalist performance in solving olympiad geometry with alphageometry2. J. Mach. Learn. Res. 26, 1--39 (2025). 题名中译: AlphaGeometry 2 解奥林匹克几何题的金牌级表现.
\bibitem{alphaproof-ref4} The mathlib Community. The Lean mathematical library. In Proc. 9th ACM SIGPLAN International Conference on Certified Programs and Proofs 367--381 (2020). 题名中译: Lean 数学库.
\bibitem{alphaproof-ref5} Dalgaard, M., Motzoi, F., Sørensen, J. J. \& Sherson, J. Global optimization of quantum dynamics with AlphaZero deep exploration. npj Quantum Inf. 6, 6 (2020). 题名中译: 用 AlphaZero 深度探索全局优化量子动力学.
\bibitem{alphaproof-ref6} Mankowitz, D. J. et al. Faster sorting algorithms discovered using deep reinforcement learning. Nature 618, 257--263 (2023). 题名中译: 用深度强化学习发现更快的排序算法.
\bibitem{alphaproof-ref7} Fawzi, A. et al. Discovering faster matrix multiplication algorithms with reinforcement learning. Nature 610, 47--53 (2022). 题名中译: 用强化学习发现更快的矩阵乘法算法.
\bibitem{alphaproof-ref8} Jaech, A. et al. OpenAI o1 system card. Preprint at \url{https://arxiv.org/abs/2412.16720} (2024). 题名中译: OpenAI o1 系统卡.
\bibitem{alphaproof-ref9} Guo, D. et al. DeepSeek-R1: incentivizing reasoning capability in LLMs via reinforcement learning. Nature 645, 633--638 (2025). 题名中译: DeepSeek-R1: 通过强化学习激励大语言模型推理能力.
\bibitem{alphaproof-ref10} Kavukcuoglu, K. Gemini 2.5: Our most intelligent AI model. Google \url{https://blog.google/technology/google-deepmind/gemini-model-thinking-updates-march-2025/} (2025). 题名中译: Gemini 2.5: 我们最智能的 AI 模型.
\bibitem{alphaproof-ref11} Petrov, I. et al. Proof or bluff? Evaluating LLMs on 2025 USA Math Olympiad. In Proc. Second AI for MATH Workshop at the 42nd International Conference on Machine Learning (eds Huang, Y. et al.) 1--15 (2025). 题名中译: 证明还是虚张声势? 在 2025 美国数学奥林匹克上评估 LLM.
\bibitem{alphaproof-ref12} Mahdavi, H. et al. Brains vs. bytes: evaluating LLM proficiency in olympiad mathematics. In Conference on Language Modeling (COLM) (eds Artzi, Y. et al.) \url{https://openreview.net/forum?id=uXR2KsA4L9} (OpenReview.net, 2025). 题名中译: 大脑对字节: 评估 LLM 的奥林匹克数学能力.
\bibitem{alphaproof-ref13} Polu, S. \& Sutskever, I. Generative language modeling for automated theorem proving. In 5th Conference on Artificial Intelligence and Theorem Proving (AITP) (eds Hales, T. C. et al.) (2020). 题名中译: 自动定理证明的生成式语言建模.
\bibitem{alphaproof-ref14} Yang, K. et al. LeanDojo: theorem proving with retrieval-augmented language models. Adv. Neural Inf. Process. Syst. 36, 21573--21612 (2023). 题名中译: LeanDojo: 用检索增强语言模型证明定理.
\bibitem{alphaproof-ref15} Vaswani, A. et al. Attention is all you need. Adv. Neural Inf. Process. Syst. 30, 5998--6008 (2017). 题名中译: 注意力就是你所需要的一切.
\bibitem{alphaproof-ref16} Li, Y. et al. Competition-level code generation with alphacode. Science 378, 1092--1097 (2022). 题名中译: AlphaCode 的竞赛级代码生成.
\bibitem{alphaproof-ref17} Lample, G. et al. Hypertree proof search for neural theorem proving. Adv. Neural Inf. Process. Syst. 35, 26337--26349 (2022). 题名中译: 神经定理证明的超树证明搜索.
\bibitem{alphaproof-ref18} Hubert, T. et al. Learning and planning in complex action spaces. In International Conference on Machine Learning 4476--4486 (PMLR, 2021). 题名中译: 复杂动作空间中的学习与规划.
\bibitem{alphaproof-ref19} Coulom, R. Computing Elo ratings of move patterns in the game of Go. ICGA J. 30, 198--208 (2007). 题名中译: 计算围棋着法模式的 Elo 评级.
\bibitem{alphaproof-ref20} Zheng, K., Han, J. M. \& Polu, S. MiniF2F: a cross-system benchmark for formal olympiad-level mathematics. In International Conference on Learning Representations (ICLR) (eds Liu, Y. et al.) \url{https://openreview.net/forum?id=9ZPegFuFTFv} (OpenReview.net, 2022). 题名中译: MiniF2F: 跨系统的形式奥林匹克级数学基准.
\bibitem{alphaproof-ref21} Tsoukalas, G. et al. PutnamBench: evaluating neural theorem-provers on the putnam mathematical competition. Adv. Neural Inf. Process. Syst. 37, 11545--11569 (2024). 题名中译: PutnamBench: 在普特南数学竞赛上评估神经证明器.
\bibitem{alphaproof-ref22} Team, G. et al. Gemini 1.5: unlocking multimodal understanding across millions of tokens of context. Preprint at \url{https://arxiv.org/abs/2403.05530} (2024). 题名中译: Gemini 1.5: 在数百万 token 上下文中释放多模态理解.
\bibitem{alphaproof-ref23} Ying, H. et al. InternLM-Math: open math large language models toward verifiable reasoning. Preprint at \url{https://arxiv.org/abs/2402.06332} (2024). 题名中译: InternLM-Math: 走向可验证推理的开放数学大语言模型.
\bibitem{alphaproof-ref24} Polu, S. et al. Formal mathematics statement curriculum learning. In International Conference on Learning Representations (eds Kim, B. et al.) \url{https://openreview.net/forum?id=-P7G-8dmSh4} (OpenReview.net, 2023). 题名中译: 形式数学命题的课程学习.
\bibitem{alphaproof-ref25} Wang, H. et al. Kimina-prover preview: towards large formal reasoning models with reinforcement learning. Preprint at \url{https://arxiv.org/abs/2504.11354} (2025). 题名中译: Kimina-Prover 预览: 通过强化学习走向大型形式推理模型.
\bibitem{alphaproof-ref26} Ren, Z. et al. Deepseek-Prover-V2: advancing formal mathematical reasoning via reinforcement learning for subgoal decomposition. Preprint at \url{https://arxiv.org/abs/2504.21801} (2025). 题名中译: DeepSeek-Prover-V2: 通过子目标分解的强化学习推进形式数学推理.
\bibitem{alphaproof-ref27} 65th IMO 2024. IMO \url{https://www.imo-official.org/year_info.aspx?year=2024} (2024). 题名中译: 第 65 届 IMO 2024.
\bibitem{alphaproof-ref28} AlphaProof and AlphaGeometry teams. AI achieves silver-medal standard solving international mathematical olympiad problems. Google DeepMind \url{https://deepmind.google/discover/blog/ai-solves-imo-problems-at-silver-medal-level/} (2024). 题名中译: AI 解国际数学奥林匹克题达到银牌标准.
\bibitem{alphaproof-ref29} Paulson, L. C. Isabelle: A Generic Theorem Prover (Springer, 1994). 题名中译: Isabelle: 通用定理证明器.
\bibitem{alphaproof-ref30} Harrison, J. HOL light: a tutorial introduction. In International Conference on Formal Methods in Computer-Aided Design 265--269 (Springer, 1996). 题名中译: HOL Light: 入门教程.
\bibitem{alphaproof-ref31} Barras, B. et al. The COQ Proof Assistant Reference Manual Version 6, 17--21 (INRIA, 1999). 题名中译: COQ 证明助手参考手册.
\bibitem{alphaproof-ref32} Gonthier, G. The four colour theorem: engineering of a formal proof. In Asian Symposium on Computer Mathematics 333--333 (Springer, 2007). 题名中译: 四色定理: 形式证明的工程.
\bibitem{alphaproof-ref33} Leroy, X. et al. CompCert: A formally verified optimizing compiler. In Embedded Real Time Software and Systems (ERTS 2016) (ed. Sifakis, J.) 1--10 (SEE, 2016). 题名中译: CompCert: 经形式验证的优化编译器.
\bibitem{alphaproof-ref34} Brown, T. et al. Language models are few-shot learners. Adv. Neural Inf. Process. Syst. 33, 1877--1901 (2020). 题名中译: 语言模型是少样本学习器.
\bibitem{alphaproof-ref35} The Claude 3 Model Family: Opus, Sonnet, Haiku, Anthropic Technical Report (Anthropic, 2024); \url{https://www.anthropic.com/claude-3-model-card} 题名中译: Claude 3 模型家族: Opus、Sonnet、Haiku.
\bibitem{alphaproof-ref36} Team, G. et al. Gemini: a family of highly capable multimodal models. Preprint at \url{https://arxiv.org/abs/2312.11805} (2023). 题名中译: Gemini: 高能力多模态模型家族.
\bibitem{alphaproof-ref37} Azerbayev, Z., Piotrowski, B. \& Avigad, J. ProofNet: A benchmark for autoformalizing and formally proving undergraduate-level mathematics problems. In Proc. NeurIPS 2022 Workshop on MATH-AI (eds Polu, S. et al.) (2022). 题名中译: ProofNet: 本科数学问题的自动形式化和形式证明基准.
\bibitem{alphaproof-ref38} Jiang, A. Q., Li, W., Han, J. M. \& Wu, Y. LISA: language models of isabelle proofs. In Proc. 6th Conference on Artificial Intelligence and Theorem Proving (eds Douglas, M. et al.) 63--65 (2021). 题名中译: LISA: Isabelle 证明的语言模型.
\bibitem{alphaproof-ref39} Bansal, K., Loos, S., Rabe, M., Szegedy, C. \& Wilcox, S. HOList: an environment for machine learning of higher order logic theorem proving. In International Conference on Machine Learning 454--463 (PMLR, 2019). 题名中译: HOList: 高阶逻辑定理证明的机器学习环境.
\bibitem{alphaproof-ref40} Han, J. M., Rute, J., Wu, Y., Ayers, E. \& Polu, S. Proof artifact co-training for theorem proving with language models. In International Conference on Learning Representations (eds Liu, Y. et al.) \url{https://openreview.net/forum?id=rpxJc9j04U} (OpenReview.net, 2022). 题名中译: 用语言模型证明定理的证明产物协同训练.
\bibitem{alphaproof-ref41} Wu, M., Norrish, M., Walder, C. \& Dezfouli, A. TacticZero: learning to prove theorems from scratch with deep reinforcement learning. Adv. Neural Inf. Process. Syst. 34, 9330--9342 (2021). 题名中译: TacticZero: 用深度强化学习从头学习证明定理.
\bibitem{alphaproof-ref42} Szegedy, C., Rabe, M. \& Michalewski, H. Retrieval-augmented proof step synthesis. In Book of Abstracts of the 6th Conference on Artificial Intelligence and Theorem Proving (AITP 2021) (eds Hales, T. et al.) 100--102 (AITP, 2021). 题名中译: 检索增强的证明步骤合成.
\bibitem{alphaproof-ref43} Welleck, S., Liu, J., Lu, X., Hajishirzi, H. \& Choi, Y. NaturalProver: grounded mathematical proof generation with language models. Adv. Neural Inf. Process. Syst. 35, 4913--4927 (2022). 题名中译: NaturalProver: 用语言模型生成有依据的数学证明.
\bibitem{alphaproof-ref44} Jiang, A. Q. et al. Thor: wielding hammers to integrate language models and automated theorem provers. Adv. Neural Inf. Process. Syst. 35, 8360--8373 (2022). 题名中译: Thor: 使用自动化工具整合语言模型与自动证明器.
\bibitem{alphaproof-ref45} First, E., Rabe, M. N., Ringer, T. \& Brun, Y. Baldur: Whole-proof generation and repair with large language models. In Proc. 31st ACM Joint European Software Engineering Conference and Symposium on the Foundations of Software Engineering (eds Chandra, S. et al.) 1229--1241 (ACM, 2023). 题名中译: Baldur: 用大语言模型生成和修复完整证明.
\bibitem{alphaproof-ref46} Jiang, A. Q. et al. Draft, sketch, and prove: Guiding formal theorem provers with informal proofs. In International Conference on Learning Representations (eds Kim, B. et al.) \url{https://openreview.net/forum?id=SMa9EAovKMC} (OpenReview.net, 2023). 题名中译: 草稿、提纲与证明: 用非形式证明指导形式证明器.
\bibitem{alphaproof-ref47} Zhao, X., Li, W. \& Kong, L. Subgoal-based demonstration learning for formal theorem proving. In Proc. 41st International Conference on Machine Learning (eds Salakhutdinov, R. et al.) Vol. 235, 60832--60865 (PMLR, 2024). 题名中译: 基于子目标示范学习的形式定理证明.
\bibitem{alphaproof-ref48} Wang, H., Wang, W., Yue, W., Sun, X. \& Wei, F. DT-Solver: Automated theorem proving with dynamic-tree sampling guided by proof-level value function. In Proc. 61st Annual Meeting of the Association for Computational Linguistics (Volume 1: Long Papers) (eds Rogers, A. et al.) 12632--12646 (2023). 题名中译: DT-Solver: 证明级价值函数指导的动态树采样自动定理证明.
\bibitem{alphaproof-ref49} Xin, H. et al. DeepSeek-Prover-V1.5: Harnessing proof assistant feedback for reinforcement learning and Monte-Carlo tree search. In Proc. Thirteenth International Conference on Learning Representations (eds Garg, A. et al.) (2025). 题名中译: DeepSeek-Prover-V1.5: 利用证明助手反馈进行强化学习和蒙特卡洛树搜索.
\bibitem{alphaproof-ref50} Wang, Q., Kaliszyk, C. \& Urban, J. First experiments with neural translation of informal to formal mathematics. In Intelligent Computer Mathematics: Proc. 11th International Conference, CICM 2018 255--270 (Springer, 2018). 题名中译: 非形式数学到形式数学的神经翻译初步实验.
\bibitem{alphaproof-ref51} Wang, Q., Brown, C., Kaliszyk, C. \& Urban, J. Exploration of neural machine translation in autoformalization of mathematics in Mizar. In Proc. 9th ACM SIGPLAN International Conference on Certified Programs and Proofs (eds Blanchette, J. \& Hrițcu, C.) 85--98 (2020). 题名中译: 在 Mizar 数学自动形式化中探索神经机器翻译.
\bibitem{alphaproof-ref52} Wu, Y. et al. Autoformalization with large language models. Adv. Neural Inf. Process. Syst. 35, 32353--32368 (2022). 题名中译: 大语言模型的自动形式化.
\bibitem{alphaproof-ref53} Agrawal, A., Gadgil, S., Goyal, N., Narayanan, A. \& Tadipatri, A. Towards a mathematics formalisation assistant using large language models. Preprint at \url{https://arxiv.org/abs/2211.07524} (2022). 题名中译: 走向利用大语言模型的数学形式化助手.
\bibitem{alphaproof-ref54} Gadgil, S., Tadipatri, A. R., Agrawal, A., Narayanan, A. \& Goyal, N. Towards automating formalisation of theorem statements using large language models. In Proc. NeurIPS 2022 Workshop on MATH-AI (eds Polu, S. et al.) 1--15 (NeurIPS, 2022). 题名中译: 走向用大语言模型自动形式化定理命题.
\bibitem{alphaproof-ref55} Jiang, A. Q., Li, W. \& Jamnik, M. Multi-language diversity benefits autoformalization. Adv. Neural Inf. Process. Syst. 37, 83600--83626 (2024). 题名中译: 多语言多样性有益于自动形式化.
\bibitem{alphaproof-ref56} Ying, H. et al. Lean workbook: a large-scale lean problem set formalized from natural language math problems. Adv. Neural Inf. Process. Syst. 37, 105848--105863 (2024). 题名中译: Lean Workbook: 从自然语言数学题形式化的大规模 Lean 问题集.
\bibitem{alphaproof-ref57} Li, Z. et al. Autoformalize mathematical statements by symbolic equivalence and semantic consistency. In Proc. 38th International Conference on Neural Information Processing Systems 1697 (Curran Associates Inc., 2024). 题名中译: 通过符号等价与语义一致性自动形式化数学命题.
\bibitem{alphaproof-ref58} Lu, J. et al. Formalalign: automated alignment evaluation for autoformalization. In International Conference on Learning Representations (eds Garg, A. et al.) \url{https://openreview.net/forum?id=B5RrIFMqbe} (OpenReview.net, 2025). 题名中译: FormalAlign: 自动形式化的自动对齐评估.
\bibitem{alphaproof-ref59} Wei, J. et al. Chain-of-thought prompting elicits reasoning in large language models. Adv. Neural Inf. Process. Syst. 35, 24824--24837 (2022). 题名中译: 思维链提示激发大语言模型推理.
\bibitem{alphaproof-ref60} Zelikman, E., Wu, Y., Mu, J. \& Goodman, N. STaR: bootstrapping reasoning with reasoning. Adv. Neural Inf. Process. Syst. 35, 15476--15488 (2022). 题名中译: STaR: 用推理自举推理.
\bibitem{alphaproof-ref61} OpenAI. OpenAI o1 system card. Preprint at \url{https://arxiv.org/abs/2412.16720} (2024). 题名中译: OpenAI o1 系统卡.
\bibitem{alphaproof-ref62} OpenAI et al. Competitive programming with large reasoning models. Preprint at \url{https://arxiv.org/abs/2502.06807} (2025). 题名中译: 大型推理模型的竞赛编程.
\bibitem{alphaproof-ref63} Jones, A. L. Scaling scaling laws with board games. Preprint at \url{https://arxiv.org/abs/2104.03113} (2021). 题名中译: 用棋类游戏扩展缩放定律.
\bibitem{alphaproof-ref64} Sun, Y. et al. Test-time training with self-supervision for generalization under distribution shifts. In International Conference on Machine Learning 9229--9248 (PMLR, 2020). 题名中译: 用自监督测试时训练应对分布偏移.
\bibitem{alphaproof-ref65} Hardt, M. \& Sun, Y. Test-time training on nearest neighbors for large language models. In Proc. Twelfth International Conference on Learning Representations (eds Chaudhuri, S. et al.) (2024). 题名中译: 大语言模型的最近邻测试时训练.
\bibitem{alphaproof-ref66} Akyürek, E. et al. The surprising effectiveness of test-time training for few-shot learning. In International Conference on Machine Learning (eds Singh, A. et al.) Vol. 267, 942--963 (PMLR, 2025). 题名中译: 测试时训练对少样本学习的惊人有效性.
\bibitem{alphaproof-ref67} Silver, D. Reinforcement Learning and Simulation-based Search in Computer Go. PhD thesis, University of Alberta (2009). 题名中译: 计算机围棋中的强化学习与基于模拟的搜索.
\bibitem{alphaproof-ref68} Zahavy, T. et al. Diversifying AI: towards creative chess with AlphaZero. Preprint at \url{https://arxiv.org/abs/2308.09175} (2023). 题名中译: 让 AI 多样化: 用 AlphaZero 走向创造性国际象棋.
\bibitem{alphaproof-ref69} Zuo, Y. et al. TTRL: test-time reinforcement learning. In Thirty-ninth Annual Conference on Neural Information Processing Systems \url{https://openreview.net/forum?id=VuVhgEiu20} (OpenReview.net, 2025). 题名中译: TTRL: 测试时强化学习.
\bibitem{alphaproof-ref70} Simonds, T. \& Yoshiyama, A. LADDER: self-improving LLMs through recursive problem decomposition. Preprint at \url{https://arxiv.org/abs/2503.00735} (2025). 题名中译: LADDER: 通过递归问题分解自我改进的 LLM.
\bibitem{alphaproof-ref71} Coquand, T. \& Paulin, C. Inductively defined types. In International Conference on Computer Logic 50--66 (Springer, 1988). 题名中译: 归纳定义类型.
\bibitem{alphaproof-ref72} Schrittwieser, J. et al. Mastering Atari, Go, chess and shogi by planning with a learned model. Nature 588, 604--609 (2020). 题名中译: 通过学习模型规划掌握 Atari、围棋、国际象棋和将棋.
\bibitem{alphaproof-ref73} Silver, D. et al. Mastering the game of Go with deep neural networks and tree search. Nature 529, 484--489 (2016). 题名中译: 用深度网络和树搜索掌握围棋.
\bibitem{alphaproof-ref74} Silver, D. et al. Mastering the game of Go without human knowledge. Nature 550, 354--359 (2017). 题名中译: 不借助人类知识掌握围棋.
\bibitem{alphaproof-ref75} Pólya, G. How to Solve It: A New Aspect of Mathematical Method (Princeton University Press, 2014). 题名中译: 怎样解题: 数学方法的新面貌.
\bibitem{alphaproof-ref76} Gonthier, G. \& Mahboubi, A. An introduction to small scale reflection in COQ. J. Formaliz. Reason. 3, 95--152 (2010). 题名中译: COQ 中小规模反射的介绍.
\bibitem{alphaproof-ref77} Hasker, R. W. \& Reddy, U. S. Generalization at higher types. In Proc. Workshop on the λProlog Programming Language (ed. Miller, D.) 257--271 (1992). 题名中译: 高阶类型上的推广.
\bibitem{alphaproof-ref78} Heras, J., Komendantskaya, E., Johansson, M. \& Maclean, E. Proof-pattern recognition and lemma discovery in ACL2. In International Conference on Logic for Programming Artificial Intelligence and Reasoning 389--406 (Springer, 2013). 题名中译: ACL2 中的证明模式识别与引理发现.
\bibitem{alphaproof-ref79} Melis, E. \& Whittle, J. Analogy in inductive theorem proving. J. Autom. Reason. 22, 117--147 (1999). 题名中译: 归纳定理证明中的类比.
\bibitem{alphaproof-ref80} Gauthier, T., Kaliszyk, C., Urban, J., Kumar, R. \& Norrish, M. TacticToe: learning to prove with tactics. J. Autom. Reason. 65, 257--286 (2021). 题名中译: TacticToe: 学习用 tactic 证明.
\end{thebibliography}
